\section{Linear constraints and integer variables}
\label{sec:constraints}

Branch-and-bound methods for mixed-integer nonlinear programming solve
constrained continuous relaxations, and the decisions they take are exact
comparisons. A node is discarded when its bound is at least the incumbent
value, a constraint is either binding or slack, and an integer assignment is
either optimal or not. For a globally strongly convex objective without
constraints, Theorem~\ref{thm:exact-upper} reduces every such comparison at
the optimizer to one $\PosSLP$ instance. This section asks what rational
linear constraints and integer variables add to that picture.

The answer has three parts.

\emph{Arbitrary polyhedra.} The set of all constraints that are tight at the
unique optimizer is a certificate of polynomial length, and it is the only
certificate that our verifier accepts. Exact comparisons therefore lie in
$\mathrm{UP}^{\PosSLP}\cap\mathrm{coUP}^{\PosSLP}$
(Theorem~\ref{thm:constraints-unambiguous}). We do not know how to find the
active set deterministically in polynomial time with a $\PosSLP$ oracle.
Section~\ref{sec:constraints-open} explains what a deterministic method would
have to supply.

\emph{Structure.} Three different kinds of structure remove the guess. If the
cubic and quartic part of the objective depends on only $k$ linear forms,
exact comparison needs no oracle and is fixed-parameter tractable in $k$
(Theorem~\ref{thm:constraints-nonlinear}). If every Taylor quadratic program
can be solved exactly with polynomially many arithmetic operations,
independently of the bit lengths of its data, then exact Newton steps give a
deterministic polynomial-time algorithm with a $\PosSLP$ oracle
(Theorem~\ref{thm:constraints-transfer}). Boxes with forest-structured or
Stieltjes Hessians and separable polynomial network flows are examples, and
for them each exact predicate reduces to a single $\PosSLP$ instance
(Corollary~\ref{cor:structured-single-sign}). If the constraint normals span a
space of small dimension $r$, a Las Vegas algorithm with a $\PosSLP$ oracle
runs in expected time $(r+1)^{O(r)}L^{C}$
(Theorem~\ref{thm:constraints-rank}).

\emph{Integer variables.} The search over integer blocks needs only ordinary
rational computation at polynomial precision. It produces a list that
contains every optimal integer block
(Theorems~\ref{thm:constraints-candidates}
and~\ref{thm:constraints-mixed-candidates}). Exact arithmetic is needed only
to compare the continuous fiber minima of the listed blocks. Knowing the
optimal value does not remove that arithmetic: one binary decision can carry
a $\PosSLP$-hard sign (Corollary~\ref{cor:constraints-binary}).

Table~\ref{tab:constraints-summary} summarizes the results. All of them assume
global strong convexity, or global strong monotonicity in the
variational-inequality versions, with a supplied rational modulus. The
established ingredients are credited where they are used: polyhedral
optimality and Farkas' lemma, essential-variable extraction, the quadratic
convergence of exact constrained Newton steps, exact box and flow quadratic
programming, linear programming with an arithmetic count independent of the
objective and right-hand-side lengths, Clarkson-type sampling
for violator spaces, and integer feasibility with integer-point separation
oracles. What this section adds to them is the exact-arithmetic analysis
that connects them: separation bounds that hold uniformly over the unknown
active face, circuit simulations whose cost does not depend on the length of
expanded rational numbers, and oracle constructions whose answers never
require an exact fiber value.

\begin{table}[t]
\centering
\small
\begin{tabularx}{\textwidth}{@{}>{\raggedright\arraybackslash}X
  >{\raggedright\arraybackslash}p{3.9cm}
  >{\raggedright\arraybackslash}p{2.9cm}
  >{\raggedright\arraybackslash}p{2.1cm}@{}}
\toprule
Additional structure & Computation & Output & Result\\
\midrule
none (arbitrary rational polyhedron)
 & $\mathrm{UP}^{\PosSLP}\cap\mathrm{coUP}^{\PosSLP}$
 & sign of $h(p)$
 & Thm.~\ref{thm:constraints-unambiguous}\\
supplied lower bound on nonzero optimal slacks
 & one $\PosSLP$ instance
 & sign of $h(p)$
 & Prop.~\ref{prop:upper-slack-gap}\\
nonlinear dimension $k$
 & ordinary, $F(k)L^{C}$
 & sign, active set, implicit $p$
 & Thm.~\ref{thm:constraints-nonlinear}\\
exact Taylor QP with polynomial arithmetic cost
 & $\mathrm P^{\PosSLP}$
 & sign, active set
 & Thm.~\ref{thm:constraints-transfer}\\
box with forest, Stieltjes, or comparison-matrix Hessians; separable flows
 & one $\PosSLP$ instance per predicate
 & sign, one active bound
 & Cor.~\ref{cor:constraints-boxes}--\ref{cor:structured-single-sign}\\
constraint rank $r$
 & Las Vegas with $\PosSLP$ oracle, expected $(r+1)^{O(r)}L^{C}$
 & sign, implicit $p$
 & Thm.~\ref{thm:constraints-rank}\\
\midrule
$t$ integer variables, unconstrained continuous fibers
 & ordinary list in $2^{O(t\log(t+1))}L^{C}$; nonadaptive $\PosSLP$ selection
 & every optimal integer block
 & Thm.~\ref{thm:constraints-candidates}, Cor.~\ref{cor:constraints-selection}\\
$t$ integer variables, mixed linear constraints
 & ordinary list in $F(t)L^{C}$
 & list containing every optimal block
 & Thm.~\ref{thm:constraints-mixed-candidates}\\
same, with nonlinear dimension $k$
 & ordinary, $F(k+t)L^{C}$
 & optimal block, value signs
 & Cor.~\ref{cor:constraints-nonlinear-mixed}\\
same, with continuous constraint rank $r$
 & Las Vegas with $\PosSLP$ oracle, expected $F(t,r)L^{C}$
 & first optimal block, implicit $p$
 & Thm.~\ref{thm:constraints-mixed-rank}\\
one binary variable, rank one, known optimal value
 & $\PosSLP$-hard
 & the optimal bit
 & Cor.~\ref{cor:constraints-binary}\\
\bottomrule
\end{tabularx}
\caption{Exact comparison at the optimizer of a globally strongly convex
explicit quartic (or the solution of a strongly monotone cubic variational
inequality) under linear constraints and integer variables. Here $F$ denotes
a computable function and $C$ an absolute constant. ``Implicit $p$'' means a
rational affine chart together with an explicit polynomial system whose unique
real solution determines $p$.}
\label{tab:constraints-summary}
\end{table}

\subsection{Setting}
\label{sec:constraints-setting}

Unless stated otherwise, $f\in\Q[x_1,\ldots,x_n]$ is an explicitly encoded
polynomial of degree at most four, and a rational $\mu>0$ is supplied with the
promise
\begin{equation}\label{eq:constraints-curvature}
 \nabla^2 f(x)\succeq\mu I\qquad(x\in\R^n).
\end{equation}
The feasible set is an explicit rational polyhedron
$P=\{x\in\R^n:Ax\le b\}$ with $A\in\Q^{m\times n}$, rows $a_i^{\mathsf T}$,
and $b\in\Q^m$. Rows may be zero, repeated, or redundant; $P$ may be unbounded
or lower-dimensional; an equality is written as two inequalities. An
\emph{observable} is an explicitly encoded $h\in\Q[x]$ of degree at most four,
and ${\bowtie}$ ranges over $\{<,\le,=,\ne,\ge,>\}$. Value and coordinate
comparisons $f(p)\bowtie\tau$ and $p_j\bowtie\tau$ with $\tau\in\Q$ are the
observables $h=f-\tau$ and $h=x_j-\tau$. We write $L\ge2$ for the total binary
length of all these data, including $\mu$, $h$, thresholds, and any supplied
certificate.

When $P\ne\varnothing$, \eqref{eq:constraints-curvature} makes $f$ coercive,
so $f$ has a unique minimizer $p$ on $P$. Its \emph{active set} is
$S(p)=\{i:a_i^{\mathsf T}p=b_i\}$, the complete set of tight rows rather than
a selected support. Emptiness of $P$ is decided by rational linear programming
before any question about $p$ is asked, and value comparisons use the
convention $\min\varnothing=+\infty$.

Several results hold for variational inequalities without a potential. Let
$T:\R^n\to\R^n$ be an explicitly encoded polynomial map of degree at most
three, with a supplied rational $\mu>0$ and the promise
\begin{equation}\label{eq:constraints-monotone}
 \bigl(T(x)-T(y)\bigr)^{\mathsf T}(x-y)\ge\mu\|x-y\|^2
 \qquad(x,y\in\R^n).
\end{equation}
A solution of $\mathrm{VI}(T,P)$ is a point $p\in P$ with
$T(p)^{\mathsf T}(x-p)\ge0$ for every $x\in P$. For $T=\nabla f$, condition
\eqref{eq:constraints-monotone} is equivalent to
\eqref{eq:constraints-curvature}, and the solutions of $\mathrm{VI}(\nabla f,P)$
are the minimizers of $f$ on $P$.

The curvature conditions are promises; nothing below recognizes them. If the
input instead carries a checked certificate, namely a positive definite full
Hessian Gram for $f$ or the Jacobian Gram of Theorem~\ref{thm:upper-monotone}
for $T$, then an invalid certificate is rejected before anything else is done.
Each statement below then becomes a statement about an ordinary language, with
a positive rational modulus derived from the certificate. The nonlinear
dimension theorem uses a weaker checkable format
(Remark~\ref{rem:constraints-psd-certificate}).

The structural parameters of this section are the nonlinear dimension $k$
(Section~\ref{sec:constraints-nonlinear}), the rank $r$ of the continuous
constraint normals (Section~\ref{sec:constraints-rank}), and the number $t$
of integer variables (Section~\ref{sec:constraints-integer}). An algorithm is
\emph{fixed-parameter tractable} (FPT) in a parameter $\kappa$ if its running
time is at most $F(\kappa)L^{C}$ for a computable function $F$ and an absolute
constant $C$. A constant that depends on the parameter never enters the
exponent of $L$.

The following geometric facts are used throughout. Their proofs are
elementary and are given in Appendix~\ref{app:constraints-geometry}. Index
sets of rows are denoted by $J$; $A_J$ and $b_J$ are the corresponding
submatrix and subvector.

\begin{lemma}[Polyhedral optimality and affine charts]
\label{lem:constraints-polyhedral}
Let $P\ne\varnothing$, let $T$ satisfy \eqref{eq:constraints-monotone}, and
for $x\in P$ write $S(x)=\{i:a_i^{\mathsf T}x=b_i\}$.
\begin{enumerate}[label=(\alph*)]
\item $\mathrm{VI}(T,P)$ has exactly one solution $p$.
\item A point $x\in P$ solves $\mathrm{VI}(T,P)$ if and only if
$-T(x)\in\operatorname{cone}\{a_i:i\in S(x)\}$. In that case there are a set
$J\subseteq S(x)$ of linearly independent rows, so $|J|\le\rank A$, and a
vector $\lambda\in\R^J_{\ge0}$ with $T(x)+A_J^{\mathsf T}\lambda=0$.
\item For every set $J$ of linearly independent rows, rational linear algebra
computes in polynomial time $\bar x\in\Q^n$ and $Z\in\Q^{n\times(n-|J|)}$
with
\[
 \{x:A_Jx=b_J\}=\{\bar x+Zy:y\in\R^{n-|J|}\},\qquad Z^{\mathsf T}Z\succeq I.
\]
Their bit lengths are bounded by one polynomial in $L$ for all $J$. The
explicit cubic map $U_J(y)=Z^{\mathsf T}T(\bar x+Zy)$ satisfies
\eqref{eq:constraints-monotone} with the same $\mu$. For $T=\nabla f$ it is
the gradient of $g_J(y)=f(\bar x+Zy)$, and $\nabla^2g_J\succeq\mu I$.
\item If $A_Jp=b_J$ and $T(p)\in\Span\{a_i:i\in J\}$, then $p=\bar x+Zy_J$,
where $y_J$ is the unique zero of $U_J$. This applies to every support $J$
given by (b) at $x=p$, and to every row basis $J$ of $A_{S(p)}$.
\end{enumerate}
\end{lemma}

When $|J|=n$, the matrix $Z$ is empty and the chart is the single rational
point $\bar x$; we then read every statement about $U_J$ and its zero as a
statement about this point. Part~(d) is what makes a guessed or computed
active set useful: it converts a constrained problem into the unconstrained
zero of an explicit strongly monotone map, to which
Theorems~\ref{thm:exact-upper} and~\ref{thm:upper-monotone} apply. Part~(c)
needs no conditioning assumption, because the identity block in $Z$ gives
$Z^{\mathsf T}Z\succeq I$.

\subsection{Arbitrary polyhedra: a unique active-set certificate}
\label{sec:constraints-general}

Lemma~\ref{lem:constraints-polyhedral}(b) suggests a certificate. Guess an
independent support $J$, compute the zero $y_J$ of $U_J$, and check primal
feasibility, nonnegativity of the multipliers
\begin{equation}\label{eq:constraints-multipliers}
 \lambda_J(y)=-(A_JA_J^{\mathsf T})^{-1}A_J\,T(\bar x+Zy),
\end{equation}
and stationarity $T(\bar x+Zy)+A_J^{\mathsf T}\lambda_J(y)=0$ at $y=y_J$.
Every tested quantity is an explicit polynomial of degree at most three
evaluated at the zero of a strongly monotone cubic map, so
Theorem~\ref{thm:upper-monotone} decides each test with $\PosSLP$. This puts
exact comparison in $\NP^{\PosSLP}\cap\mathrm{coNP}^{\PosSLP}$, but the
certificate is not unique: with degenerate multipliers many supports pass. The
complete active set $S(p)$ is unique. Verifying it, however, requires showing
that $-T(p)$ lies in the cone of the active normals without knowing which
multipliers to use. That is a linear program whose objective is the algebraic
vector $T(p)$.

\begin{theorem}[Unambiguous certificates]
\label{thm:constraints-unambiguous}
Let $T$ satisfy \eqref{eq:constraints-monotone}, let $P=\{x:Ax\le b\}$ be an
explicit rational polyhedron, let $h$ be an observable, and let
${\bowtie}\in\{<,\le,=,\ne,\ge,>\}$. Consider the promise problem whose
yes-instances satisfy $P\ne\varnothing$ and $h(p)\bowtie0$, where $p$ is the
unique solution of $\mathrm{VI}(T,P)$. It belongs to
$\mathrm{UP}^{\PosSLP}\cap\mathrm{coUP}^{\PosSLP}$. In particular, for an
explicit quartic $f$ satisfying \eqref{eq:constraints-curvature}, every exact
value, coordinate, and observable comparison at the minimizer of $f$ on $P$
belongs to $\mathrm{UP}^{\PosSLP}\cap\mathrm{coUP}^{\PosSLP}$. With a checked
Hessian or Jacobian Gram in place of the bare modulus, the same holds for the
ordinary languages in which invalid certificates are rejected.
\end{theorem}

\begin{proof}[Proof idea]
The machine first rejects invalid certificates and decides $P\ne\varnothing$
by linear programming. It then guesses one bit per row, that is, a set
$S\subseteq\{1,\ldots,m\}$. Deterministically it chooses a row basis $J$ of
$A_S$, the chart of Lemma~\ref{lem:constraints-polyhedral}(c), and the zero
$p_S=\bar x+Zy_J$ of $U_J$. It accepts the optimality part of the certificate
exactly when two conditions hold: the slack $b_i-a_i^{\mathsf T}p_S$ is zero
for $i\in S$ and positive for $i\notin S$; and
\begin{equation}\label{eq:constraints-cone-test}
 T(p_S)^{\mathsf T}v\ge0\quad\text{for every }v\in Q_S,
 \qquad Q_S=\{v:A_Sv\le0,\ -1\le v_j\le1\ (j\le n)\}.
\end{equation}
Every slack is an affine observable at the zero of $U_J$. The test
\eqref{eq:constraints-cone-test} is the delicate step. Every vertex $v$ of
$Q_S$ has coordinates of polynomial bit length, so all vertex observables
$v^{\mathsf T}T(\bar x+Zy)$ have one common encoding bound. One shared Newton
circuit $\hat y$ for $y_J$ (Lemma~\ref{lem:newton-circuit}) therefore
approximates all of them to within one eighth of their common separation
bound. A linear-programming algorithm whose arithmetic count depends only on
the encoded constraint matrix, simulated on
rational circuits with $\PosSLP$ comparisons, minimizes the rational circuit
objective $T(\bar x+Z\hat y)^{\mathsf T}v$ over $Q_S$. Its optimal vertex is
recovered as printed rational numbers, and one exact comparison of its true
observable decides \eqref{eq:constraints-cone-test}. A passing set $S$ forces
$p_S=p$ and $S=S(p)$, and $S(p)$ passes. Thus exactly one guess passes, and
the machine finally tests $h(p_S)\bowtie0$, or its complement for the
$\mathrm{coUP}$ side. Appendix~\ref{app:constraints-unambiguous} gives the
complete proof, including the circuit-objective linear program
(Lemma~\ref{lem:constraints-circuit-lp}).
\end{proof}

\begin{remark}[What the theorem does and does not give]
\label{rem:constraints-unambiguous}
Guessing an independent support instead of the full mask gives the simpler
$\NP^{\PosSLP}\cap\mathrm{coNP}^{\PosSLP}$ bound described above
(Lemma~\ref{lem:constraints-support-test}). Enumerating all supports is
deterministic but takes $\sum_{s\le n}\binom ms$ steps. By
\citet{AllenderEtAl2009}, $\PosSLP$ lies in the counting hierarchy, hence so do
all predicates of Theorem~\ref{thm:constraints-unambiguous}. The theorem does
not reduce the constrained problem to one $\PosSLP$ instance, and the adaptive
compilation of Theorem~\ref{thm:posslp-closure}(b) does not apply, because the
guess of the active set is not removed. Nor does ordinary approximation
find the active set in general: nothing in the separation bound prevents a
nonzero slack at $p$ from being doubly exponentially small in $L$, and an
approximation to polynomially many bits cannot distinguish such a slack from
zero.
\end{remark}

If the input also supplies a margin that separates the active rows from the
inactive ones, ordinary approximation finds the active set, and the oracle is
needed only for one final comparison. Proposition~\ref{prop:upper-slack-gap}
states this: when every nonzero optimal slack is at least a supplied rational
$\delta>0$, an exactly feasible rational approximation of the minimizer
determines every slack to within $\delta/8$, thresholding at $\delta/2$
returns $S(p)$, and
Theorem~\ref{thm:exact-upper}, applied on the chart of a row basis of
$A_{S(p)}$, yields one $\PosSLP$ instance per relation. No lower bound on
nonzero multipliers is needed. The promise is substantial and is not verified,
and the separation bound for algebraic numbers does not supply a margin of
polynomial bit length.

\subsection{Few nonlinear directions}
\label{sec:constraints-nonlinear}

Many models are quadratic except for a few higher-degree terms in a few
linear combinations of the variables, for example a strongly convex quadratic
plus a quartic penalty $(a^{\mathsf T}x)^4$, or a quartic risk term in a few
factors. The constraints of such a model may be numerous and of full rank.
The next theorem shows that the algebraic difficulty of the optimizer then
lives in few dimensions, and that exact comparison needs no oracle.

\begin{definition}[Nonlinear dimension]
\label{def:constraints-nonlinear-dimension}
Write $f=f_{\le2}+f_3+f_4$, where $f_j$ is homogeneous of degree $j$ and
$f_{\le2}$ has degree at most two. The \emph{nonlinear kernel} of $f$ is
\[
 K_f=\{v\in\R^n:D_vf_3\equiv0\text{ and }D_vf_4\equiv0\},
\]
where $D_vg=\sum_iv_i\,\partial_ig$ and $\equiv$ denotes equality of
polynomials. The \emph{nonlinear dimension} of $f$ is $k=n-\dim K_f$.
\end{definition}

\begin{lemma}[Essential directions]
\label{lem:constraints-essential}
\begin{enumerate}[label=(\alph*)]
\item Let $W$ be the rational matrix with rows indexed by the monomials
$x^\beta$ of degrees two and three, whose entry $W_{\beta,i}$ is the
coefficient of $x^\beta$ in $\partial_i(f_3+f_4)$. Then $K_f=\ker W$. Hence
$k=\rank W$, and rational matrices $U\in\Q^{k\times n}$ with $\ker U=K_f$ and
$V\in\Q^{n\times k}$ with $UV=I_k$ are computable in polynomial time.
\item $f_3(x)+f_4(x)=\theta(Ux)$ for all $x$, where
$\theta(u)=f_3(Vu)+f_4(Vu)$.
\item The nonlinear dimension does not change under an invertible rational
affine change of variables.
\end{enumerate}
\end{lemma}

Part~(b) is the essential-variable construction; for homogeneous forms it is
\citet[Proposition~1]{Carlini2006}. Here it is applied separately to the cubic
and quartic parts, and part~(c) shows that the resulting parameter is
intrinsic even though translation mixes these parts.

\begin{theorem}[Exact comparison in few nonlinear directions]
\label{thm:constraints-nonlinear}
Let $f$ satisfy \eqref{eq:constraints-curvature} and have nonlinear
dimension $k$, let $P=\{x:Ax\le b\}$ be an explicit rational polyhedron, and
let $h$ be an observable. There are a function $F(k)=2^{O(k)}$ and an
absolute constant $C$ such that a deterministic ordinary algorithm, using no
oracle, runs in time $F(k)L^{C}$ and either reports $P=\varnothing$ or
\begin{enumerate}[label=(\roman*)]
\item decides whether $h(p)$ is negative, zero, or positive;
\item computes the active set $S(p)$; and
\item computes $\bar c\in\Q^n$ and $D\in\Q^{n\times s}$ of full column rank
$s\le k$ such that $p=\bar c+Du_*$, where $u_*$ is the unique zero of the
gradient of the explicit quartic $\varphi(u)=f(\bar c+Du)$ on $\R^s$, and
$\nabla^2\varphi\succeq\mu D^{\mathsf T}D\succeq\nu I$ with the rational
modulus $\nu=\mu\det(D^{\mathsf T}D)/\tr(D^{\mathsf T}D)^{s-1}>0$ when $s\ge1$.
\end{enumerate}
The data in (iii) have bit length at most $T(L)$ for a polynomial $T$ that
does not depend on $k$. The number of constraints, their rank, the ambient
dimension, and the rank of the quadratic part are unrestricted.
\end{theorem}

\begin{proof}[Proof idea]
The proof uses the unknown active set only to establish a bound. Let $J$ be
a row basis of $A_{S(p)}$ and $\bar x+Zy$ its chart. By
Lemma~\ref{lem:constraints-polyhedral}(d), $p$ is the unconstrained minimizer
of $f(\bar x+Zy)$. Split $y$ along $\ker(UZ)$ and a complement. On the kernel
directions the cubic and quartic parts are constant, so $f$ is a strongly
convex quadratic there and can be minimized in closed form. What remains is
a strongly convex quartic $\varphi$ in $s=\rank(UZ)\le k$ variables, with
data of bit length $T(L)$ for every possible $J$. The formula
$\exists u\,(\nabla\varphi(u)=0\wedge z=h(\bar c+Du))$ has $s$ quantified
variables and defines the singleton $\{h(p)\}$. The coefficient bound in
one-block quantifier elimination is linear in the input coefficient length,
with a factor depending only on the number of quantified variables
(Lemma~\ref{lem:separation}). Hence $h(p)\ne0$ implies
$|h(p)|\ge2^{-G(k)T(L)}$ with $G(k)=2\cdot4^{c_0k}$, where $c_0$ is the
absolute constant of that lemma. An exactly feasible rational
point with objective error $2^{-O(G(k)T(L))}$ is within distance
$2^{-G(k)T(L)}/(8B)$ of $p$, where $B$ bounds $\|\nabla h\|$, and is
computable in ordinary time $F(k)L^{C}$ (Lemma~\ref{lem:convex-value}).
Comparing its rational observable value with $\pm2^{-G(k)T(L)-1}$ decides
the sign exactly. The slacks give $S(p)$, and then the elimination above is
carried out on the now known set. Appendix~\ref{app:constraints-nonlinear}
gives the details.
\end{proof}

The algorithm prints $G(k)T(L)$ accuracy bits. This number is FPT in $k$,
but for unconstrained objectives of nonlinear dimension $n$ it is
exponential in $n$, and Theorem~\ref{thm:exact-upper} is then the better
bound. The
two results measure different sources of difficulty.

\begin{example}\label{ex:constraints-rank-one-quartic}
Let $f(x)=\tfrac12x^{\mathsf T}Qx+c^{\mathsf T}x+(a^{\mathsf T}x)^4$ with
rational $Q\succeq\mu I$, $a\ne0$, and $c$. Then $K_f=a^\perp$ and $k=1$.
Theorem~\ref{thm:constraints-nonlinear} decides every exact comparison at the
minimizer over an arbitrary rational polyhedron in ordinary polynomial time,
although $Q$ may be dense, the constraint matrix may have rank $n$, and the
polyhedron may have exponentially many faces. The optimizer is an explicit
affine function of the unique real root of one univariate cubic equation.
\end{example}

\begin{remark}[A checkable curvature certificate]
\label{rem:constraints-psd-certificate}
A rational positive semidefinite matrix $G_0$ with
$v^{\mathsf T}(\nabla^2f(x)-\mu I)v=w^{\mathsf T}G_0w$, where
$w=(v,x\otimes v)$, certifies \eqref{eq:constraints-curvature}. Coefficient
matching and rational semidefiniteness testing check it in polynomial time.
A positive definite full Hessian Gram is the wrong format here: if $k<n$, it
cannot exist. Indeed, for $0\ne e\in K_f$ the Hessian $\nabla^2f(se)$ does not
depend on $s$, whereas a positive definite Gram would give
$v^{\mathsf T}\nabla^2f(se)v\ge c(1+s^2\|e\|^2)\|v\|^2$ for some $c>0$. We do
not claim that every globally strongly convex quartic has a certificate of
either kind.
\end{remark}

\begin{remark}[Related parameters]
\label{rem:constraints-nonlinear-prior}
Low-rank optimization usually bounds the rank of the whole objective
\citep{MittalSchulz2013} and asks for approximation. \citet{KannanRademacher2009}
treat a convex function plus a polynomial perturbation in $k$ variables, again
with an approximation guarantee whose cost grows with the reciprocal
accuracy. Here only the cubic and quartic part is restricted, the quadratic
part and the constraints are arbitrary, the output is exact including zero,
and the exponent of $L$ does not depend on $k$. The ingredients that make this
possible are strong convexity, which turns objective accuracy into point
accuracy and eliminates the quadratic directions, and the linear dependence of
the elimination bound on coefficient length.
\end{remark}

\subsection{Exact constrained Newton steps}
\label{sec:constraints-newton}

For an unconstrained problem, a Newton step is a linear solve, so exact
Newton steps can be carried out on shared rational circuits; this is the
mechanism behind Theorem~\ref{thm:exact-upper}. With constraints, a Newton
step is a quadratic program over $P$. Its convergence is classical. The
difficulty is arithmetic: the Taylor data of the $j$-th step can have bit
length exponential in $j$, so a quadratic-programming algorithm whose running
time is polynomial in the bit length of its data is of no use. What is needed is an
algorithm whose number of arithmetic operations does not depend on the bit
length of the Hessian and linear cost.

\begin{definition}[Arithmetic Taylor-QP solver]
\label{def:constraints-qp-solver}
Let $\mathcal C$ be a class of instances $(f,P)$ with $f$ explicit,
satisfying \eqref{eq:constraints-curvature}, and $P\ne\varnothing$. An
\emph{arithmetic Taylor-QP solver} for $\mathcal C$ is a deterministic
algorithm with the following properties. It receives the explicit data of an
instance in $\mathcal C$ and rational $H\in\Q^{n\times n}$, $c\in\Q^n$ such
that $H=\nabla^2f(x)$ and $c=\nabla f(x)-\nabla^2f(x)x$ for some $x\in P$. It
returns the exact minimizer of $\tfrac12y^{\mathsf T}Hy+c^{\mathsf T}y$ over
$P$. Its operations on numbers are addition, subtraction, multiplication,
division by a nonzero number, and comparison, applied to the entries of $H$
and $c$, to earlier results, and to rational constants computed from the
explicit data in time polynomial in $L$. The choice of each operation is
determined by the outcomes of earlier comparisons through a control that runs
in time polynomial in $L$, and the number of operations is at most a
polynomial in $L$ for all admissible $(H,c)$, independently of their bit
lengths.
\end{definition}

A dependence on logarithms of numerical bounds that are themselves of
polynomial bit length is allowed; a dependence on the bit lengths of $H$ or
$c$ is not. A quadratic-programming algorithm that is polynomial only in the
bit length of its data does not satisfy the definition.

\begin{lemma}[Constrained Newton step]
\label{lem:constraints-newton-step}
Let $f$ be twice continuously differentiable with $\nabla^2f\succeq\mu I$, let
$P$ be a nonempty closed convex set with minimizer $p$ of $f$, and let $x\in P$.
Let $\Lambda$ be a Lipschitz constant of $\nabla^2f$, in operator norm, on a
convex set containing $x$ and $p$. Then the minimizer
\[
 N(x)=\argmin_{y\in P}\Bigl\{\nabla f(x)^{\mathsf T}(y-x)
 +\tfrac12(y-x)^{\mathsf T}\nabla^2f(x)(y-x)\Bigr\}
\]
satisfies $\|N(x)-p\|\le\frac{\Lambda}{2\mu}\|x-p\|^2$.
\end{lemma}

\begin{proof}
Write $y=N(x)$ and $H=\nabla^2f(x)$. The first-order conditions of the two
convex problems give
$(\nabla f(x)+H(y-x))^{\mathsf T}(p-y)\ge0$ and
$\nabla f(p)^{\mathsf T}(y-p)\ge0$. Add them and put
$\rho=\nabla f(p)-\nabla f(x)-H(p-x)$. Then
\[
 \mu\|y-p\|^2\le(y-p)^{\mathsf T}H(y-p)\le\rho^{\mathsf T}(y-p)
 \le\|\rho\|\,\|y-p\|.
\]
Taylor's formula gives
$\rho=\int_0^1\bigl(\nabla^2f(x+\vartheta(p-x))-H\bigr)(p-x)\,d\vartheta$, so
$\|\rho\|\le\frac\Lambda2\|x-p\|^2$.
\end{proof}

This is the exact proximal Newton estimate of
\citet[Theorem~3.4]{LeeSunSaunders2014}, with the indicator function of $P$ as
the nonsmooth term. We record the proof because it shows a point that matters
here: the estimate never identifies the optimal face. It holds when the Taylor
model and the true problem have different active sets, and when optimal
multipliers vanish.

\begin{theorem}[Transfer from arithmetic QP solvers]
\label{thm:constraints-transfer}
Fix $D\ge2$. Let $\mathcal C$ be a class of instances $(f,P)$ with $f$
explicit of degree at most $D$ satisfying \eqref{eq:constraints-curvature}
and $P$ a nonempty explicit rational polyhedron, and suppose that
$\mathcal C$ has an arithmetic Taylor-QP solver. For every explicit observable
$h$ of degree at most $D$, each predicate $h(p)\bowtie0$ is decided in
deterministic polynomial time with a $\PosSLP$ oracle, and $S(p)$ is computed
in the same way. If $f(p)<\tau$ for a rational $\tau$, the computation also
yields a point $x\in P$ with $f(x)<\tau$, given by a shared rational circuit
of polynomial size.
\end{theorem}

\begin{proof}[Proof idea]
An ordinary convex-optimization call gives an exactly feasible rational warm
start $x_0$ with $\|x_0-p\|\le1/(2\kappa)$, where $\kappa=\max\{1,\Lambda/(2\mu)\}$
and $\Lambda$ is an explicit Hessian-Lipschitz bound on a ball that contains
all iterates. Lemma~\ref{lem:constraints-newton-step} gives
$\|x_j-p\|\le\kappa^{-1}2^{-2^j}$ for the exact iterates $x_{j+1}=N(x_j)$.
Each iterate is computed by the solver, with every number held as a shared
numerator-denominator circuit and every comparison answered by $\PosSLP$.
The Karush--Kuhn--Tucker system of the original instance, extended by
$z=h(x)$, defines the singleton $\{h(p)\}$, so Lemma~\ref{lem:separation}
gives $h(p)=0$ or $|h(p)|\ge2^{-2^{e_D(L)}}$ for a polynomial $e_D$ that
depends only on $D$. After polynomially many steps a final comparison of
$h(x_j)$ with $\pm2^{-2^{e_D(L)}-1}$ decides the predicate.
The comparisons inside the solver depend on earlier oracle answers, so the
procedure is a Turing reduction. Appendix~\ref{app:constraints-newton} gives
the complete argument.
\end{proof}

The theorem isolates the arithmetic interface; it does not assert that
arbitrary convex quadratic programs have such solvers. We now give two classes
that do.

For a symmetric matrix $H$, its \emph{comparison matrix} $M(H)$ has
$M(H)_{ii}=H_{ii}$ and $M(H)_{ij}=-|H_{ij}|$ for $i\ne j$.

\begin{corollary}[Structured boxes]
\label{cor:constraints-boxes}
Let $P=\{x:\ell\le x\le u\}$ with $\ell_i\in\Q\cup\{-\infty\}$ and
$u_i\in\Q\cup\{+\infty\}$, and let $f$ be explicit of fixed degree $D$ and
satisfy \eqref{eq:constraints-curvature}. Suppose that one of the following
holds for every $x\in P$:
\begin{enumerate}[label=(\roman*)]
\item every pair $i\ne j$ with $\partial_i\partial_jf(x)\ne0$ is an edge of a
supplied forest on $\{1,\ldots,n\}$;
\item $\partial_i\partial_jf(x)\le0$ for all $i\ne j$;
\item $M(\nabla^2f(x))$ is positive definite.
\end{enumerate}
Then for every explicit observable $h$ of degree at most $D$, the predicates
$h(p)\bowtie0$ are in $\mathrm P^{\PosSLP}$, and so is the set of active
bounds. The number of active bounds and the dimension are unrestricted, and
zero multipliers are allowed.
\end{corollary}

Conditions (i) and (ii) imply (iii). For (ii), $M(H)=H$. For (i), choosing
signs along each tree gives a diagonal sign matrix $\Sigma$ with
$M(H)=\Sigma H\Sigma$. The proof (Appendix~\ref{app:constraints-boxes})
first truncates infinite bounds without changing $p$
(Lemma~\ref{lem:constraints-truncation}). It then shows that every positive
definite $H$ with $M(H)\succ0$ admits a rational vector $v\ge\mathbf 1$,
computable by rational linear algebra and sign tests, such that
$H_{SS}^{-1}v_S\ge0$ for every nonempty $S$
(Lemma~\ref{lem:constraints-admissible}). With such a vector, a parametric
principal pivoting method in the style of \citet[Algorithm~I]{PangHan2023}
solves the box quadratic program with at most $2n$ pivots and $O(n^4)$
arithmetic operations (Lemma~\ref{lem:constraints-box-qp}). We include
complete proofs of both lemmas so that the operation count is explicit.
Condition~(i) can be checked from the coefficients of $f$ when the box is
full-dimensional. Conditions (ii) and (iii) are promises.

For any graph $G$, rational $w_{ij}\ge0$, and rational $c$,
\[
 f(x)=\frac\mu2\|x\|^2+\sum_{ij\in E(G)}w_{ij}(x_i-x_j)^4+c^{\mathsf T}x
\]
has Hessian $\mu I$ plus a graph Laplacian with nonnegative weights
$12w_{ij}(x_i-x_j)^2$, so it satisfies (ii) on every box. The graph may be
dense, the box constraints have rank $n$, and the nonlinear dimension can be
$n-1$. For example, if $G$ is a star with center $1$ and positive weights,
the forms $x_1-x_j$ are linearly independent, so their cubes are linearly
independent and $K_f$ is spanned by the all-ones vector. Neither
Theorem~\ref{thm:constraints-nonlinear} nor Theorem~\ref{thm:constraints-rank}
then gives a polynomial bound. Replacing each difference by
$x_i+x_{i+1}$ along a path gives a tridiagonal Hessian with positive
off-diagonal entries, covered by (i) but not by (ii).

\begin{corollary}[Separable network flows]
\label{cor:constraints-flows}
Let $G=(V,E)$ be a directed graph with incidence matrix $A_G$, and let
$P=\{x\in\R^E:A_Gx=b,\ \ell\le x\le u\}$ with rational $b$ and with
$\ell_e\in\Q\cup\{-\infty\}$, $u_e\in\Q\cup\{+\infty\}$. Let
$f(x)=\sum_{e\in E}f_e(x_e)$ with explicit $f_e\in\Q[s]$ of fixed degree $D$
and $f_e''\ge\mu$ on $\R$. Then for every explicit observable $h$ of degree at
most $D$, the predicates $h(p)\bowtie0$ are in $\mathrm P^{\PosSLP}$, and so is
the set of active bounds. The graph, the dimension, and the constraint rank
are unrestricted; parallel arcs, redundant balance equations, and fixed arcs
are allowed.
\end{corollary}

The Taylor model of a separable objective is a separable strictly convex
quadratic flow problem with arc costs
$\tfrac12f_e''(x_e)y_e^2+\bigl(f_e'(x_e)-f_e''(x_e)x_e\bigr)y_e$. The strongly
polynomial algorithm of \citet[Theorem~20]{Vegh2016} solves it exactly with a
number of arithmetic operations and comparisons that depends only on the size
of the graph. Appendix~\ref{app:constraints-flows} states the precise
interface and proves that the artificial arc costs used by such algorithms can
be chosen from the circuit data. For $D\le4$ the curvature promise is
checkable: $f_e''-\mu$ is a univariate quadratic. The observable $h$ may
couple many arcs; only the objective must be separable.

The deterministic oracle algorithms of
Corollaries~\ref{cor:constraints-boxes} and~\ref{cor:constraints-flows} make
adaptive queries. The compilation theorem for deterministic $\PosSLP$
computations removes the adaptivity.

\begin{corollary}[One sign instance for structured classes]
\label{cor:structured-single-sign}
Under the hypotheses of Theorem~\ref{thm:constraints-transfer}, and in
particular of Corollary~\ref{cor:constraints-boxes} or
Corollary~\ref{cor:constraints-flows}, there is a deterministic
polynomial-time algorithm that maps each instance, observable $h$, and relation
${\bowtie}$ to one integer arithmetic circuit $\Phi$ such that, on every
promised instance, $\Phi>0$ if and only if $h(p)\bowtie0$. The same holds for
each predicate ``bound $i$ is active at $p$'' and for every polynomial-size
Boolean combination of such predicates.
\end{corollary}

\begin{proof}
The algorithm of Theorem~\ref{thm:constraints-transfer} is a deterministic
oracle Turing machine whose queries are integer circuits. Its running time
and the lengths of its queries are bounded by one polynomial in $L$ that does
not depend on the expanded values of the circuits: the warm start is an
ordinary computation, the number of Newton steps and of solver operations per
step is polynomial in $L$, and each query is a circuit that the machine has
written. A clock and syntax checks make the machine total on all inputs
without changing it on promised ones.
Theorem~\ref{thm:posslp-closure}(b) compiles such a machine into one circuit.
Appendix~\ref{app:constraints-single-sign} checks these hypotheses in detail.
\end{proof}

The corollary concerns one requested predicate at a time. The set of active
bounds is a vector of bits; it is obtained from a nonadaptive list of compiled
instances, one per bound, not from one circuit. The structured classes are not
claimed to be $\PosSLP$-hard, and no ordinary polynomial-time algorithm is
claimed. Theorem~\ref{thm:constraints-unambiguous} and the Las Vegas
Theorem~\ref{thm:constraints-rank} are not covered: the first guesses its
certificate, and the second uses randomness with an expected rather than a
worst-case time bound. For parameterized compositions, for example the
candidate lists of Section~\ref{sec:constraints-integer} with structured
fibers, the same compilation gives a many-one reduction whose construction
time is FPT in the parameter, not polynomial when the parameter grows.

\subsection{Few constraint directions}
\label{sec:constraints-rank}

A third parameter is the rank $r=\rank A$ of the constraint normals. A
polyhedron can have very many constraints whose normals lie in a space of
small dimension, for example many inequalities on a few linear forms of the
variables. By Lemma~\ref{lem:constraints-polyhedral}(b) the solution is then
determined by at most $r$ active rows, but enumerating the candidate supports
takes $m^{O(r)}$ steps, with $r$ in the exponent of the input length.
Random sampling in the style of Clarkson removes $r$ from that exponent.

\begin{theorem}[Constraint rank]
\label{thm:constraints-rank}
Let $T$ satisfy \eqref{eq:constraints-monotone}, let $P=\{x:Ax\le b\}$ be an
explicit rational polyhedron with $r=\rank A$, and let $h$ be an observable.
There is a Las Vegas algorithm with a $\PosSLP$ oracle that decides whether
$P=\varnothing$ and otherwise decides any predicate $h(p)\bowtie0$ at the
unique solution $p$ of $\mathrm{VI}(T,P)$. Its expected running time is
$(r+1)^{O(r)}L^{C}$ for an absolute constant $C$, every oracle query has
length polynomial in $L$, and every answer it returns is correct. It also
returns a set $J$ of at most $r$ linearly independent rows such that, in the
notation of Lemma~\ref{lem:constraints-polyhedral}, $p=\bar x+Zy_J$ with
$y_J$ the unique zero of the explicit map $U_J$. For $T=\nabla f$ this covers
exact value, coordinate, and observable comparisons for strongly convex
quartic optimization.
\end{theorem}

\begin{proof}[Proof idea]
After checking that $P$ is nonempty, every subset $G$ of the rows defines a
nonempty polyhedron $P_G$ with a unique solution $p_G$ of
$\mathrm{VI}(T,P_G)$. The map $V(G)=\{i:a_i^{\mathsf T}p_G>b_i\}$ defines a
violator space in the sense of \citet{GaertnerMatousekRuestSkovron2008}.
Locality follows from uniqueness, so no potential function is needed.
Applying Lemma~\ref{lem:constraints-polyhedral}(b) to an inclusion-minimal
basis shows that every basis consists of linearly independent rows, so the
combinatorial dimension is at most $r$. The sampling algorithms need only
violation tests whose first argument has at most $r$ elements. Such a test
enumerates the at most $2^r$ independent subsets of its argument and checks
each with the support test of Lemma~\ref{lem:constraints-support-test}, using
Theorem~\ref{thm:upper-monotone}. Clarkson's algorithms, as analyzed for
violator spaces by \citet{GaertnerMatousekRuestSkovron2008}, use an expected
$O(rm+r^{O(r)})$ such tests. Appendix~\ref{app:constraints-rank} gives the
complete proof, including the binary implementation of the weighted sampling.
\end{proof}

The parameter is the rank of all constraint normals, not the size of the
final active support; the sampling analysis bounds the bases of all
subproblems. A box has rank $n$, so the theorem does not make box-constrained
problems easy. The oracle is essential even for $r=0$: unconstrained exact
comparison is already $\PosSLP$-hard (Theorem~\ref{thm:quartic-complete}).
Randomness affects only the running time. We do not claim a deterministic
algorithm with this parameter dependence, and the compilation of
Corollary~\ref{cor:structured-single-sign} does not apply to this algorithm.

\subsection{Integer variables}
\label{sec:constraints-integer}

We now add $t$ integer variables. Consider
\begin{equation}\label{eq:constraints-mixed-problem}
 \min\bigl\{f(z,y):z\in\Z^t,\ y\in\R^n,\ (z,y)\in P\bigr\},
\end{equation}
where $f\in\Q[z,y]$ is explicit of degree at most four with a supplied
rational $\mu>0$ and the promise $\nabla^2f(z,y)\succeq\mu I_{t+n}$ on all of
$\R^{t+n}$. We treat two kinds of feasible sets: $P=Q\times\R^n$ with an
explicit rational polyhedron $Q\subseteq\R^t$ (unconstrained continuous
fibers), and $P=\{(z,y):Az+By\le d\}$ with explicit rational $A$, $B$, $d$
(mixed linear constraints). An \emph{optimal integer block} is the $z$-part of
an optimal solution. For $z$ in the projection of $P$ onto the $z$-coordinates,
write $g(z)=\min\{f(z,y):(z,y)\in P\}$ for its fiber minimum; $z$ is
\emph{feasible} if it is an integer point of this projection. When
\eqref{eq:constraints-mixed-problem} is feasible, coercivity on the closed set
$(\Z^t\times\R^n)\cap P$ gives an optimal solution, each fiber has a unique
minimizer, and $g$ is $\mu$-strongly convex on the projection of $P$.

The fiber values $g(z)$ are algebraic numbers, typically irrational. An
integer optimization method that compares them, or bisects on the objective,
needs exact arithmetic at every step. Our search avoids exact values
altogether. At each queried integer point it computes, at polynomial
precision, a rational inequality that is violated by the query and satisfied
by every other integer point whose fiber value is no larger. The following
lemma is the source of all such cuts.

\begin{lemma}[Integer cut]
\label{lem:constraints-integer-cut}
Let $z\in\Z^t$ be feasible and let $c_z\in\Q^t$ satisfy
\begin{equation}\label{eq:constraints-integer-cut}
 g(w)-g(z)\ge c_z^{\mathsf T}(w-z)+\tfrac\mu4\|w-z\|^2
\end{equation}
for every feasible $w\in\Z^t\setminus\{z\}$. If $c_z=0$, then $z$ is the
unique optimal integer block. Otherwise the inequality
$c_z^{\mathsf T}(x-z)\le-\mu/4$ is violated by $z$ and satisfied by every
feasible integer $w\ne z$ with $g(w)\le g(z)$, in particular by every optimal
integer block other than $z$.
\end{lemma}

\begin{proof}
If $c_z=0$, \eqref{eq:constraints-integer-cut} gives $g(w)>g(z)$ for every
feasible integer $w\ne z$. Otherwise $g(w)\le g(z)$ and $\|w-z\|\ge1$ give
$c_z^{\mathsf T}(w-z)\le-\frac\mu4\|w-z\|^2\le-\frac\mu4$.
\end{proof}

Such a cut may remove points of the continuous sublevel set
$\{x:g(x)\le g(z)\}$, including nonintegral points of smaller value. It is
therefore not a separation oracle for that convex set. We use integer
feasibility algorithms whose oracle is queried only at integer points and
must separate a fixed closed convex set from each rejected query
(Theorem~\ref{thm:constraints-integer-query}). The proofs fix the target set
after the fact: the empty set for the running-time bound, and a single
omitted optimal block for completeness.

\begin{theorem}[Candidate lists with unconstrained fibers]
\label{thm:constraints-candidates}
Let $P=Q\times\R^n$ with an explicit rational polyhedron $Q\subseteq\R^t$,
possibly unbounded or lower-dimensional. A deterministic ordinary algorithm,
using no oracle, either certifies $Q\cap\Z^t=\varnothing$ or outputs a finite
list $\mathcal Z\subseteq Q\cap\Z^t$ that contains every optimal integer
block. Its running time and the total encoding length of $\mathcal Z$ are at
most $2^{O(t\log(t+1))}L^{C}$, with $C$ independent of $t$ and $n$. There are
at most $2^t$ optimal integer blocks, and this bound is attained.
\end{theorem}

\begin{proof}[Proof idea]
Fixed-dimensional integer linear feasibility gives a feasible block, and
strong convexity gives a radius $R$, of bit length $2^{O(t)}L^{O(1)}$, that
contains every optimal block. The fiber minimum $g$ is differentiable with
$\nabla g(z)=\nabla_zf(z,y(z))$, where $y(z)$ is the fiber minimizer. An
ordinary approximation of $y(z)$ at polynomial precision gives $c_z$ with
$\|c_z-\nabla g(z)\|\le\mu/4$, and strong convexity then gives
\eqref{eq:constraints-integer-cut}. The integer-query feasibility algorithm is
run with these cuts and is never told that a point is feasible. Every point it
queries in $Q\cap[-R,R]^t$ is recorded. Viewed against the empty target set,
the run is valid, which bounds its length. If an optimal block $w$ were never
queried, the same run would be a valid run for the target $\{w\}$, and the
algorithm would wrongly report that $\{w\}$ has no integer point. The parity
bound follows because two optimal blocks with the same parity have an integer
midpoint with strictly smaller value. Appendix~\ref{app:constraints-lists}
gives the details.
\end{proof}

The list separates discrete search from exact arithmetic. The search uses
ordinary rational computation at polynomial precision, and its length is FPT
in $t$ with an exponent of $L$ that is independent of both $t$ and $n$. Exact
arithmetic is needed only to compare the finitely many fiber minima of the
list. Two fiber minima are compared through one observable: $g(z)-g(z')$ is
the value of $f(z,y)-f(z',y')$ at the minimizer of the separable strongly
convex quartic $f(z,y)+f(z',y')$, so no comparison of two independently
represented algebraic numbers is needed.

Fixed-dimensional integer minimization from first-order information is
classical; the closest predecessor is the method of
\citet{OertelWagnerWeismantel2014}, whose oracle contract differs from the
integer-point separation interface used here. The integer feasibility theorem
and the idea of avoiding bisection on the objective by keeping visited
candidates are imported \citep{AriHildebrand2026,HildebrandGoess2024,Basu2023Integers}.
Those results do not by themselves optimize an algebraic fiber value. What is
added here is an implementation of the integer-point interface for the fiber
minimum from rational data of polynomial precision, the proof that every
optimal block is listed, and the extension to mixed linear constraints without
a constraint qualification.

\begin{corollary}[Exact selection by nonadaptive queries]
\label{cor:constraints-selection}
In the setting of Theorem~\ref{thm:constraints-candidates}, let $\tau\in\Q$.
\begin{enumerate}[label=(\alph*)]
\item Each predicate $\min\bowtie\tau$ for problem
\eqref{eq:constraints-mixed-problem} is a Boolean combination of the
predicates $g(z)\bowtie'\tau$ with $z\in\mathcal Z$ and
${\bowtie'}\in\{<,\le\}$. Each of these reduces to one $\PosSLP$ instance by
Theorem~\ref{thm:exact-upper}, and all instances are built from $\mathcal Z$
before any oracle answer is known. By Theorem~\ref{thm:posslp-closure}(c), the
predicate therefore has a many-one reduction to a single $\PosSLP$ instance,
computable in time $2^{O(t\log(t+1))}L^{C}$.
\item The set of all optimal integer blocks is computed in time
$2^{O(t\log(t+1))}L^{C}$ with at most $|\mathcal Z|^2$ nonadaptive $\PosSLP$
queries. For an optimal block $z_*$, the continuous part of the optimal
solution is the unique real solution of $\nabla_yf(z_*,y)=0$.
\end{enumerate}
In particular, for each fixed $t$, exact optimization of
\eqref{eq:constraints-mixed-problem} with unconstrained fibers is in
$\mathrm P^{\PosSLP}$, and each predicate $\min\bowtie\tau$ has a
polynomial-time many-one reduction to $\PosSLP$.
\end{corollary}

The construction time in (a) is FPT; it is polynomial only for bounded $t$.
The oracle cannot be removed unless $\PosSLP\in\mathrm P$, even for fixed $t$:
the case $t=0$ is unconstrained exact comparison, which is $\PosSLP$-hard
(Theorem~\ref{thm:quartic-complete}), and new integer variables can be added
without changing the optimal value. With the padding of
Lemma~\ref{lem:constraints-padding} the instances even keep a positive
definite full Hessian Gram, so the hardness holds in the certified format for
every fixed $t\ge1$. Corollary~\ref{cor:constraints-selection}
returns the complete optimal set; a sequential tournament with $|\mathcal Z|-1$
adaptive comparisons returns one optimal block.

\begin{theorem}[Candidate lists with mixed linear constraints]
\label{thm:constraints-mixed-candidates}
Let $P=\{(z,y):Az+By\le d\}$ with explicit rational data. A deterministic
ordinary algorithm, using no oracle, either certifies that
\eqref{eq:constraints-mixed-problem} is infeasible or outputs a finite list of
feasible integer blocks that contains every optimal integer block. Its running
time and the total length of the list are at most $F(t)L^{C}$ for a computable
$F$ and an absolute constant $C$. The polyhedron may be unbounded or
lower-dimensional, with redundant or degenerate constraints; no Slater
condition and no bound on Lagrange multipliers is assumed. There are at most
$2^t$ optimal integer blocks.
\end{theorem}

\begin{proof}[Proof idea]
Fibers may now be infeasible, and $g$ need not be differentiable. At an
integer query with an infeasible fiber, a rational Farkas certificate gives a
cut that is valid for the projection of $P$. At a query $z$ with a feasible
fiber, let $\hat y$ be an exactly feasible rational fiber point and
$\lambda\ge0$ a rational multiplier vector. Joint strong convexity and the
multiplier inequality give
\[
 g(w)-g(z)\ge c_z^{\mathsf T}(w-z)+\tfrac\mu2\|w-z\|^2-E,\qquad
 c_z=\nabla_zf(z,\hat y)+A^{\mathsf T}\lambda,
\]
for every feasible $w$, where
$E=\lambda^{\mathsf T}\sigma+\|\nabla_yf(z,\hat y)+B^{\mathsf T}\lambda\|^2/(2\mu)$
and $\sigma=d-Az-B\hat y\ge0$. If $E\le\mu/4$, this is
\eqref{eq:constraints-integer-cut}. To make $E$ small without a constraint
qualification, choose $\hat y$ close to the fiber minimizer by an ordinary
convex approximation. A true optimal multiplier $\lambda_*$ exists by
Lemma~\ref{lem:constraints-polyhedral}(b), and complementary slackness gives
the identity
$\lambda_*^{\mathsf T}\sigma=\nabla_yf(z,y_z)^{\mathsf T}(\hat y-y_z)$, where
$y_z$ is the fiber minimizer. This bounds the multiplier term without bounding
$\|\lambda_*\|$. The best rational multiplier then comes from an exact convex
quadratic program in $\lambda$. Appendix~\ref{app:constraints-mixed} gives the
details.
\end{proof}

The sharper parameter function $2^{O(t\log(t+1))}$ of
Theorem~\ref{thm:constraints-candidates} is not claimed here; it depends on the
initial mixed-integer feasibility step. Exact selection among the listed
blocks is a separate problem, because the fibers are now constrained. Each of
the constrained results of this section gives a selection method, with its own
computational model.

\begin{corollary}[Mixed selection in few nonlinear directions]
\label{cor:constraints-nonlinear-mixed}
In the setting of Theorem~\ref{thm:constraints-mixed-candidates}, let $k$ be
the nonlinear dimension of $f$ as a polynomial on $\R^{t+n}$. A deterministic
ordinary algorithm, using no oracle, decides each predicate $\min\bowtie\tau$,
returns every optimal integer block, and represents the corresponding
continuous optimizers as in Theorem~\ref{thm:constraints-nonlinear}, in time
$F(k+t)L^{C}$.
\end{corollary}

\begin{theorem}[Mixed constraints and constraint rank]
\label{thm:constraints-mixed-rank}
In the setting of Theorem~\ref{thm:constraints-mixed-candidates}, let
$r=\rank B$. A Las Vegas algorithm with a $\PosSLP$ oracle either detects
infeasibility or returns the lexicographically smallest optimal integer block
$z_*$ together with a representation $y_*=\bar y+Zu_*$ of the optimal
continuous part, where $u_*$ is the unique zero of an explicit strongly
monotone cubic map. It also decides each predicate $\min\bowtie\tau$, and can
return every optimal integer block. Its expected running time is
$F(t,r)L^{C}$ for a computable $F$ and an absolute constant $C$, and every
returned answer is correct.
\end{theorem}

The fibers of a pair of candidates are compared on their product, whose
constraint matrix $\diag(B,B)$ has rank $2r$, by
Theorem~\ref{thm:constraints-rank}; each fiber of $f$ has the curvature bound
$\mu$, and the nonlinear dimension of a product of fibers is at most $2k$.
Appendix~\ref{app:constraints-mixed} proves both compositions. The rank of
$A$ does not enter: its effect is absorbed by the ordinary list.

Without any structure on the fibers, the list still reduces the mixed problem
to the general constrained bound.

\begin{corollary}[Mixed unambiguous certificates]
\label{cor:constraints-mixed-unambiguous}
In the setting of Theorem~\ref{thm:constraints-mixed-candidates}, the language
of instances that are feasible with $\min\bowtie\tau$ is decided by an
unambiguous nondeterministic machine with a $\PosSLP$ oracle, and so is its
complement, both in time $F(t)L^{C}$. In particular, for each fixed $t$ these
predicates are in $\mathrm{UP}^{\PosSLP}\cap\mathrm{coUP}^{\PosSLP}$.
\end{corollary}

The machine computes the list, guesses one active set for each listed fiber,
and verifies all of them as in Theorem~\ref{thm:constraints-unambiguous}.
Exactly one guess passes for each fiber. The fiber minima are then the
unconstrained minima of explicit quartics on charts, which are compared as
above (Appendix~\ref{app:constraints-mixed}).

Knowing the optimal value does not make the integer decision easy. In the
following problem there are two candidate integer blocks, the optimal value is
known in advance, the optimizer is rational and bounded, and the continuous
constraint normals have rank one; nevertheless, the optimal bit is
$\PosSLP$-hard.

\begin{corollary}[A hard binary decision at a known optimum]
\label{cor:constraints-binary}
There is a polynomial-time map from integer straight-line programs to
instances
\begin{equation}\label{eq:constraints-binary}
 \min\Bigl\{H(x,z)=F(x)+\bigl(z-\tfrac12\bigr)^2:\
 x\in\R^N,\ z\in\Z,\ 0\le z\le1,\ z-1\le x_j\le z\Bigr\},
\end{equation}
where $F$ is a rational quartic with $N+1$ rational quadratic square factors
and a positive definite rational full Hessian Gram, such that: the optimal
value is $1/4$; the optimal solution $(x_*,z_*)$ is unique, with
$x_*\in\Q^N\cap[-1,1]^N$; $\nabla^2H\succeq\frac32I$; the continuous
constraint normals have rank one; and $z_*=1$ exactly when the program's
output is positive. Adding $\epsilon z(z-1)\|x\|^2+\frac14z^2(z-1)^2$ with a
suitable rational $\epsilon>0$ changes no feasible objective value and gives
the objective a positive definite rational full Hessian Gram on all $N+1$
variables.
\end{corollary}

Subtracting $\frac14$ from $H$ gives an equivalent instance with known
optimal value zero; the optimal solution, the Hessian, and the full Hessian
Gram matrices do not change.

The proof (Appendix~\ref{app:constraints-binary}) starts from
Theorem~\ref{thm:rational-optimizer}: the designated coordinate $p_j$ of the
rational optimizer of $F$ is nonzero, and the constraints admit $x=p$ exactly
for the bit $z=\mathbf 1[p_j>0]$. A polynomial-time algorithm that returns the
optimal integer block of \eqref{eq:constraints-binary} would therefore decide
$\PosSLP$. The gap between the two fiber values lies between
$\frac34p_j^2$ and $\frac{\Lambda_F}2p_j^2$, where $\Lambda_F$ bounds
$\|\nabla^2F\|$ on $[-1,1]^N$. No inverse-polynomial lower bound on $|p_j|$ is
assumed, and Lemma~\ref{lem:separation} guarantees only a doubly
exponentially small one. That separation bound gives no guarantee that fiber
values approximated to polynomially many bits certify the bit. The Las Vegas oracle algorithm of
Theorem~\ref{thm:constraints-mixed-rank} cannot be replaced by an ordinary
polynomial-time algorithm, even for $t=r=1$, unless $\PosSLP\in\mathrm P$.

\subsection{The open case: arbitrary polyhedra and unrestricted structure}
\label{sec:constraints-open}

For an arbitrary rational polyhedron, unrestricted nonlinear dimension, and
unrestricted constraint rank, the strongest bound proved here is
Theorem~\ref{thm:constraints-unambiguous}. Exact comparison is
$\PosSLP$-hard already without constraints. Whether it is in
$\mathrm P^{\PosSLP}$ is open; by Theorem~\ref{thm:posslp-closure}(b) this is the
same as asking for a polynomial-time many-one reduction to one $\PosSLP$
instance. We record what the methods of this section need and why three
natural routes fall short. These are limitations of particular proofs, not
lower bounds.

\emph{Exact quadratic programming with circuit Hessians.}
Lemma~\ref{lem:constraints-newton-step} does not depend on the polyhedron.
Theorem~\ref{thm:constraints-transfer} would therefore settle the general case
if every Taylor quadratic program over an arbitrary rational polyhedron could
be solved exactly with a number of arithmetic operations that is polynomial in
the encoding of $A$ and $b$ and independent of the encoding of the Hessian and
the linear cost. The classical polynomial-time algorithm for convex quadratic
programming \citep{KozlovTarasovKhachiyan1980} is polynomial in the bit length
of all data, which grows exponentially along the Newton iteration.
\citet{GranotSkorinKapov1990} remove the dependence on the linear cost and the
right-hand side, but not on the Hessian. The linear program with a circuit
objective in the proof of Theorem~\ref{thm:constraints-unambiguous} handles a
varying linear objective over a fixed polytope, not a varying Hessian. We know
of no algorithm with the required interface for general polyhedra and
Hessians; the sufficient condition is not claimed to be necessary.

\emph{Penalties.} One could try to remove the constraints by a penalty and
apply Theorem~\ref{thm:exact-upper}. For $\min_{x\ge0}\frac12(x+1)^2$, with
minimizer $0$, the squared-hinge penalty
$\frac12(x+1)^2+\frac\lambda2\max\{0,-x\}^2$ has minimizer
$-1/(1+\lambda)$. Reaching distance $\gamma$ requires
$\lambda\ge\gamma^{-1}-1$. At the separation scale
$\gamma=2^{-2^{\poly(L)}}$ allowed by Lemma~\ref{lem:separation}, the penalty coefficient
needs exponentially many printed bits, and the penalty is piecewise quadratic
rather than polynomial. A short circuit for $\lambda$ does not help, because
Theorem~\ref{thm:exact-upper} accepts explicitly encoded polynomials and
starts from an ordinary warm start of polynomial bit length.

\emph{Accuracy-driven iteration.} An interior-point or ellipsoid method with
iteration count polynomial in $\log(1/\varepsilon)$ needs exponentially many
iterations at $\varepsilon=2^{-2^{\poly(L)}}$. Storing $\varepsilon$ as a circuit
does not shorten the iteration. A rapid refinement phase with a polynomial
number of steps, which exact Newton steps provide when the subproblems can be
solved, is the missing ingredient.

The positive results of this section supply such refinement or avoid it in
different ways: a small nonlinear dimension makes the separation bound singly
exponential, so that ordinary approximation suffices; structured quadratic
programs make exact Newton steps arithmetic; low constraint rank confines the
combinatorial search to small subproblems. None of them covers a general
strongly convex quartic over a general polyhedron.
